% Editable manuscript fragment. Load the packages/macros listed in main.tex.
\section{Experiments}
\label{sec:exp}

We evaluate \methodname\ on two complementary benchmark tasks and one
retrospective case study: reaction-conditioned enzyme screening on EnzymeMap,
bidirectional enzyme--reaction retrieval on ReactZyme, and literature-supported
catalyst recovery for the conversion of D-fructose to D-tagatose. Following the
task-based organization of FGW-CLIP \citep{zhou2026fgwclip}, each benchmark is
described in terms of its data, comparators, evaluation metrics, and results.
We then analyze controlled changes to the training objective and inference
settings using the recorded ablation experiments.

\paragraph{Common experimental design.}
\methodname\ denotes one architecture and hyperparameter recipe fitted
\emph{separately} on each target training partition: the three ReactZyme splits
and the EnzymeMap split. The trainable F3 retrieval layers are freshly
initialized for every target; target-trained ReactZyme weights are not reused
to fit EnzymeMap. Frozen protein, molecular, reaction, and SLEEC feature
extractors remain pretrained resources. Phase 1 uses the anchor-balanced,
decoupled multi-positive objective and the two attention regularizers described
in the companion methods section. Phase 2 freezes F3 and fits the existing
residual heads with full-graph uniform-positive cross-entropy and an
identity-preservation penalty. The training-reaction dictionary is constructed
from the target training graph only. Neither ensembling nor hubness correction
is used.

The primary recipe retains phase-1 epoch 10 and phase-2 update 100 with seed
42. Its inference settings are dictionary weight $\alpha=0.40$, protein fusion
multiplier $\nu=3$, and residual multiplier $c=0.5$. The last parameter scales
the residual update; it is not a norm-clipping threshold. \vfourBio\ adds
relative biological supervision only to phase 2, with weight 1 for each of
the EC, cofactor, and coarse transformation families, margin 0.1, and a fixed
family divisor of three. This variant shares the original F3 checkpoints and
adds no inference-time parameters. Detailed optimization settings appear in
Appendix~\ref{app:exp-implementation}.

Matching downstream training associations is necessary for our comparisons,
but does not match every pretrained representation or annotation resource.
We therefore distinguish comparisons between complete methods from controlled
component ablations. The final recipe was reached through repeated exploratory
benchmark inspection; fixing its final checkpoints does not make the overall
development process a single-use held-out test. Reported differences are point
estimates unless a paired-seed analysis is explicitly identified.

\subsection{Enzyme screening on EnzymeMap}
\label{sec:exp-screening}
\subsubsection{Datasets and screening protocols}

We use the EnzymeMap release and reaction-rule split distributed with
CLIPZyme \citep{heid2023enzymemap,mikhael2024clipzyme}. The processed dataset
contains 46,356 enzyme--reaction association rows, 16,776 unique reactions,
12,749 enzymes, 2,841 EC identifiers, and 394 reaction rules. The released
training, validation, and test partitions contain 34,427, 7,287, and 4,642
rows, respectively. The nominal 80/10/10 allocation is made over reaction
rules; it does not imply these exact proportions of association rows. We
retain the released partitions and physical reactant/product sides rather
than resampling the split.

The screening library combines EnzymeMap, BRENDA release 2022\_2, and UniProt
release 2022\_01, with a maximum sequence length of 650 residues. Its 261,907
candidate identifiers and their ordering are retained. Query construction
follows the released CLIPZyme analysis notebook: the 4,642 test associations
contain 1,551 distinct reaction strings, of which 30 also occur in training.
Removing those reactions leaves 1,521 unique test queries. Known positive
candidate identifiers are collected from the released raw EnzymeMap records
and intersected with the screening library, yielding 4,544 query--candidate
positive labels. These labels are used for scoring, not to extend the training
association graph.

We report two screening settings. \emph{Full pool} ranks all 261,907 candidates
for each of the 1,521 queries. \emph{Training-ID-excluded pool} removes every
candidate identifier observed in the training associations, leaving 252,113
candidates and 1,337 queries with at least one remaining positive. This second
setting follows the additional generalization analysis in FGW-CLIP. It tests
exclusion of training identifiers, not an independently established
sequence-disjoint or homology-disjoint split. The existing ablation figures
refer to these settings as Table 1 and Table 2, respectively; those labels are
protocol names inherited from the source comparison, not manuscript table
cross-references.

\subsubsection{Baselines}

We compare with the published CLIPZyme encoder variants and FGW-CLIP
\citep{mikhael2024clipzyme,zhou2026fgwclip}. We also evaluate the released
CLIPZyme checkpoint on the same query lists, candidate order, positive labels,
and screening metric implementation used for \methodname. Its downstream
training data are the corresponding released EnzymeMap training partition;
it is a checkpoint reproduction, not a new training run. Published FGW-CLIP
numbers remain explicitly marked as reported values: no official checkpoint
or verified reproducible implementation was available in the recorded source
audit. Thus the CLIPZyme comparison has local score-level evidence, whereas
the FGW-CLIP comparison is against its reported point estimates.

\subsubsection{Evaluation metrics}

We measure early retrieval with BEDROC at $a=85$ and $a=20$, and enrichment
factors over the top 5\% and 10\% of each ranked library. Higher values are
better. The screening tables show BEDROC as a percentage; the ablation
figures and saved numerical files use its unscaled value. Enrichment factors
are fold enrichments, not percentages. Metrics are calculated per reaction
and averaged over eligible queries with equal query weight.

For a query with $N$ candidates, $n$ positives, and binary ranked labels
$y_{(j)}$, the released enrichment calculation is
\begin{equation}
  \operatorname{EF}_{f}
  =\frac{\sum_{j=1}^{\lfloor fN\rfloor}y_{(j)}}{f n},
  \qquad f\in\{0.05,0.10\}.
  \label{eq:exp-ef}
\end{equation}
We preserve the notebook's rank origin, sorting convention, and finite-library
normalization for BEDROC; Appendix~\ref{app:exp-metrics} gives the executable
definition. MRR is not the screening endpoint or the criterion for interpreting
these experiments.

\subsubsection{Results}

Tables~\ref{tab:exp-screen-full} and~\ref{tab:exp-screen-excluded} present all
four screening metrics. In the full pool, \methodname\ achieves BEDROC$_{85}$
of 57.2337\%, BEDROC$_{20}$ of 75.5433\%, EF5 of 16.8590, and EF10 of 8.9587.
Its BEDROC$_{85}$ exceeds the reproduced CLIPZyme value by 12.5394 percentage
points and the reported FGW-CLIP value by 8.5737 percentage points. After
excluding training enzyme identifiers, BEDROC$_{85}$ remains 52.8734\%, with
EF5 of 16.4186 and EF10 of 8.7931. These results support improved early
recovery under both recorded screening protocols.

\begin{table}[htbp]
\centering\small
\caption{EnzymeMap screening against the full library: 1,521 queries and
261,907 candidate identifiers. BEDROC is expressed as a percentage. Published
rows come from FGW-CLIP Table 1; locally measured rows use the shared released
CLIPZyme notebook protocol. All \methodname\ rows use the same target-specific
training associations.}
\label{tab:exp-screen-full}
\resizebox{\linewidth}{!}{% Generated from evidence_snapshot.json by prepare_tables.py.
\begin{tabular}{lrrrr}
\toprule
Method & BEDROC$_{85}$ (\%) & BEDROC$_{20}$ (\%) & EF5 & EF10 \\
\midrule
CLIPZyme (ESM), reported & 36.91 & 53.04 & 11.93 & 6.84 \\
CLIPZyme (CGR), reported & 38.91 & 57.58 & 13.16 & 7.73 \\
CLIPZyme, reported & 44.69 & 62.98 & 14.09 & 8.06 \\
FGW-CLIP, reported & 48.66 & 66.69 & 14.91 & 8.18 \\
\midrule
CLIPZyme, released checkpoint & 44.6942 & 62.9817 & 14.0851 & 8.0600 \\
\methodname & 57.2337 & 75.5433 & 16.8590 & 8.9587 \\
\vfourBio & 57.2096 & 75.5322 & 16.8609 & 8.9565 \\
\bottomrule
\end{tabular}
}
\end{table}

\begin{table}[htbp]
\centering\small
\caption{EnzymeMap screening after excluding training enzyme identifiers:
1,337 eligible queries and 252,113 candidate identifiers. Published values
are from FGW-CLIP Table 2. The strongest published EF10 comparator in this
setting is CLIPZyme (7.81), not FGW-CLIP (7.61).}
\label{tab:exp-screen-excluded}
\resizebox{\linewidth}{!}{% Generated from evidence_snapshot.json by prepare_tables.py.
\begin{tabular}{lrrrr}
\toprule
Method & BEDROC$_{85}$ (\%) & BEDROC$_{20}$ (\%) & EF5 & EF10 \\
\midrule
CLIPZyme, reported & 39.13 & 58.86 & 13.40 & 7.81 \\
FGW-CLIP, reported & 45.14 & 61.43 & 13.57 & 7.61 \\
\midrule
CLIPZyme, released checkpoint & 39.1328 & 58.8619 & 13.3967 & 7.8081 \\
\methodname & 52.8734 & 72.6403 & 16.4186 & 8.7931 \\
\vfourBio & 52.8464 & 72.6275 & 16.4357 & 8.8031 \\
\bottomrule
\end{tabular}
}
\end{table}

Adding biological supervision in phase 2 preserves the margins over those
reported comparators, but does not improve all metrics over \methodname.
BEDROC$_{85}$ decreases slightly in both pools; EF changes are small and
mixed. We therefore attribute the screening result to the evaluated full
recipe, without treating the biological objective as an established cause of
the gain. The controlled label-removal and shuffle experiments below address
that distinction directly.
\FloatBarrier

\subsection{Bidirectional enzyme--reaction retrieval on ReactZyme}
\label{sec:exp-reactzyme}
\subsubsection{Datasets and splits}

ReactZyme links SwissProt proteins to Rhea reactions and provides temporal,
enzyme-similarity, and reaction-similarity evaluation partitions
\citep{hua2024reactzyme}. The released resource is described as containing
178,463 association rows, 178,327 proteins, and 7,726 reaction identifiers.
We use the existing materialized \texttt{reactzyme\_paper} partitions shared
by \methodname\ and the public-baseline campaign. Their exact retained row
counts appear in Table~\ref{tab:exp-reactzyme-data}; these are the counts
actually fitted and scored, rather than nominal totals copied from the paper.
The materialized partitions contain 178,459 retained association rows per
split. Published raw split sizes and retained distinct associations should
not be interchanged.

Reaction similarity separates held-out reaction identifiers; enzyme similarity
tests a sequence-based partition; and the temporal split separates annotations
around the released 31 December 2010 cutoff. We preserve the supplied
train/validation assignment, approximately a 90/10 subdivision of the
non-test partition, and do not draw a new validation set for each comparator.
The materialized reaction inputs are flattened participant sets. Their
self-reaction encoding must not be interpreted as recovered physical
reactant/product roles. EnzymeMap, in contrast, supplies those roles.

\begin{table}[htbp]
\centering\small
\caption{Audited ReactZyme materialization. Train, validation, and test columns
count retained association rows. The last two columns give the complete test
reaction and enzyme catalogs: their roles as query and candidate sets swap
between retrieval directions. Every locally retrained comparator uses these
same catalogs and their recorded ordering.}
\label{tab:exp-reactzyme-data}
\resizebox{\linewidth}{!}{% Generated from evidence_snapshot.json by prepare_tables.py.
\begin{tabular}{lrrrrr}
\toprule
Split & Train & Validation & Test & Reactions & Enzymes \\
\midrule
Reaction similarity & 147,393 & 16,377 & 14,689 & 386 & 14,688 \\
Enzyme similarity & 152,748 & 16,972 & 8,739 & 1,573 & 8,734 \\
Time & 149,554 & 16,618 & 12,287 & 2,634 & 12,277 \\
\bottomrule
\end{tabular}
}
\end{table}

\subsubsection{Baselines and comparison scope}

We separate two comparison tracks. The \emph{published-reference track}
includes FGW-CLIP's EC Mode and EC Max variants and the two main TIGER
variants, ESM2Text and ProtT3 \citep{zhou2026fgwclip,zhang2026tiger}. TIGER's
main Reaction-Sim \expEtoR\ MRR is 0.518 for ESM2Text; the 0.543 MLP ablation
and 0.632 SwissProt-text condition are different experimental settings and
are not relabeled as that main model.

The \emph{matched-data retraining track} fits the four released ReactZyme
families---MLP, contrastive MLP, Transformer, and Bi-RNN---with both ESM2 and
SaProt protein features and four molecular feature choices (MAT-2D, MAT-3D,
UniMol-2D, and UniMol-3D). This defines 96 family/feature/split configurations.
An additional 24 configurations test a corrected contrastive-label convention
and remain separate from the unchanged released contrastive model. We also
include the completed original Horizyn participant-set adapter
\citep{rocks2026horizyn} and EnzGFM backbone controls. All completed
configurations in the frozen evidence snapshot are retained in
Appendix~\ref{app:exp-public}; test performance is not used to select a
representative feature combination.

The local baseline training uses the same downstream associations as
\methodname, fresh retrieval weights, validation-only checkpoint selection,
and a shared test evaluator. Material adaptations remain explicit: the
released classifiers use reproducible sampled negatives because their
original negative files are unavailable; SaProt follows the released ESM
fallback for missing structural strings; Horizyn consumes a participant-set
reaction adapter; and the cached EnzGFM preprocessing differs from its native
paper convention. Their objectives, training budgets, and backbones are not
identical. Consequently these are matched-data \emph{method} comparisons,
not isolated tests of architecture alone. Pending CREEP, structural CLIPZyme,
EnzymeCAGE, and VenusRXN experiments have no completed result asserted here.

\subsubsection{Evaluation metrics}

For each direction, we rank the complete corresponding test candidate catalog.
Let $P(q)$ be the positive candidates for query $q$, and
$\operatorname{rank}_{q}(p)$ their one-based ranks. Our primary metric is the
released benchmark implementation's all-positive, query-averaged MRR:
\begin{equation}
 \operatorname{MRR}_{\mathrm{all}}
 =\frac{1}{|Q|}\sum_{q\in Q}
   \frac{1}{|P(q)|}\sum_{p\in P(q)}
   \frac{1}{\operatorname{rank}_{q}(p)}.
 \label{eq:exp-mrr}
\end{equation}
We additionally retain first-positive MRR and Hit@$k$ for $k=1,5,10$ as
diagnostics (Appendix~\ref{app:exp-metrics}); they are not interchangeable
with Eq.~\eqref{eq:exp-mrr}. Scores are ranked in FP32 with stable
candidate-index tie breaking in the local shared evaluator.

Some published descriptions, including FGW-CLIP's metric formula, instead
describe first-positive MRR. Without its predictions or runnable evaluator,
we cannot independently reconcile that wording with its reported tables.
Published MRR columns are therefore reproduced \emph{as reported} and kept
separate from the common-evaluator retraining table. Local comparisons retain
one metric definition and identical candidate axes throughout.

\subsubsection{Results}

Table~\ref{tab:exp-reactzyme-literature} places \methodname\ and \vfourBio\
alongside the published reference values. Original \methodname\ records
\expRtoE/\expEtoR\ MRR of 0.400381/0.534366 on Reaction-Sim,
0.666249/0.971355 on Enzyme-Sim, and 0.538285/0.780669 on Time. These exceed
the displayed main-method literature point estimates, subject to the metric
and resource qualifications above. The biology variant produces small
\expRtoE\ changes and larger, mixed \expEtoR\ changes; numerical near ties
qualify the latter, as examined in Section~\ref{sec:exp-ties}.

\begin{table}[htbp]
\centering\small
\caption{ReactZyme literature-reference comparison. The FGW-CLIP rows are
transcribed from Tables 3, 7, and 8 of arXiv:2512.08508v2; TIGER rows are its
main Table 1 values. Published MRR is retained as reported, while our rows use
Eq.~\eqref{eq:exp-mrr}. This table is not a local reproduction of TIGER or
FGW-CLIP, nor proof of equivalent pretraining resources or metric code.}
\label{tab:exp-reactzyme-literature}
\resizebox{\linewidth}{!}{% Generated from evidence_snapshot.json by prepare_tables.py.
\begin{tabular}{lrrrrrr}
\toprule
Method & \multicolumn{2}{c}{Reaction similarity} & \multicolumn{2}{c}{Enzyme similarity} & \multicolumn{2}{c}{Time} \\
\midrule
 & R$\to$E & E$\to$R & R$\to$E & E$\to$R & R$\to$E & E$\to$R \\
\midrule
FGW-CLIP, EC Mode & 0.3181 & 0.3722 & 0.5253 & 0.8683 & 0.3394 & 0.5222 \\
FGW-CLIP, EC Max & 0.3113 & 0.3804 & 0.5300 & 0.8581 & 0.3392 & 0.5229 \\
TIGER (ESM2Text) & 0.3190 & 0.5180 & 0.5920 & 0.9560 & 0.3660 & 0.6900 \\
TIGER (ProtT3) & 0.3370 & 0.4720 & 0.5790 & 0.9400 & 0.3720 & 0.6830 \\
\midrule
\methodname & 0.400381 & 0.534366 & 0.666249 & 0.971355 & 0.538285 & 0.780669 \\
\vfourBio & 0.400412 & 0.534777 & 0.666234 & 0.967373 & 0.538317 & 0.800614 \\
\bottomrule
\end{tabular}
}
\end{table}

\begin{table}[htbp]
\centering\small
\caption{Completed matched-data ReactZyme comparisons under the shared
all-positive MRR evaluator. For the four released families and corrected-loss
sensitivity, this compact table shows the same ESM2/MAT-2D features throughout;
the EnzGFM control also uses MAT-2D. This is a fixed presentation slice, not
the best feature pair chosen on test. All completed feature combinations are
listed in Appendix~\ref{app:exp-public}. Training budgets differ across methods.}
\label{tab:exp-reactzyme-retrained}
\resizebox{\linewidth}{!}{% Generated from evidence_snapshot.json by prepare_tables.py.
\begin{tabular}{lrrrrrr}
\toprule
Method & \multicolumn{2}{c}{Reaction similarity} & \multicolumn{2}{c}{Enzyme similarity} & \multicolumn{2}{c}{Time} \\
\midrule
 & R$\to$E & E$\to$R & R$\to$E & E$\to$R & R$\to$E & E$\to$R \\
\midrule
MLP & 0.105690 & 0.218712 & 0.215499 & 0.707968 & 0.104607 & 0.284913 \\
Contrastive MLP & 0.040553 & 0.239766 & 0.119823 & 0.666660 & 0.040501 & 0.255020 \\
Transformer & 0.122008 & 0.231432 & 0.216317 & 0.725992 & 0.141887 & 0.394288 \\
Bi-RNN & 0.151047 & 0.283470 & 0.273579 & 0.796985 & 0.127391 & 0.375185 \\
Corrected contrastive MLP & 0.063832 & 0.243126 & 0.104661 & 0.631571 & 0.041036 & 0.241179 \\
EnzGFM backbone control & 0.135390 & 0.296464 & 0.286377 & 0.798907 & 0.152498 & 0.386741 \\
Horizyn participant-set adapter & 0.431815 & 0.534799 & 0.699701 & 0.976990 & 0.598203 & 0.840621 \\
\midrule
\methodname & 0.400381 & 0.534366 & 0.666249 & 0.971355 & 0.538285 & 0.780669 \\
\vfourBio & 0.400412 & 0.534777 & 0.666234 & 0.967373 & 0.538317 & 0.800614 \\
\bottomrule
\end{tabular}
}
\end{table}

The newer matched-data results give a more restrictive conclusion
(Table~\ref{tab:exp-reactzyme-retrained}). The original Horizyn
participant-set adapter obtains higher MRR than original \methodname\ in
all six split/direction cells. Its Reaction-Sim \expRtoE\ MRR is 0.431815
versus 0.400381, and its Time \expRtoE\ MRR is 0.598203 versus 0.538285.
Horizyn also exceeds \vfourBio\ in all six cells, although the
Reaction-Sim \expEtoR\ margin is small (0.534799 versus 0.534777). Thus the completed
local evidence does not support a claim that either V4 variant outperforms
every retrained competitor. Different training budgets prevent assigning
this difference to architecture alone. The full results are retained rather
than excluding the stronger parent-model comparator.
\FloatBarrier

\subsection{Case1: retrospective recovery in a small literature panel}
\label{sec:exp-case1}
\subsubsection{Reaction and candidate construction}

Case1 concerns the direct C4 epimerization of free D-fructose to D-tagatose,
as investigated in the literature summarized by the source workbook
\citep{caseS01,caseS02,caseS03,caseS04,caseS05,caseS06,caseS16,caseS17}.
The workbook's small homolog panel contains \textbf{145 entries}. One entry,
H017 (KoT4E M6), lacks a fully specified recoverable sequence. The recorded
evaluation therefore ranks \textbf{144 resolved entries}, representing
\textbf{123 unique amino-acid sequences}. We report both the original-entry
and deduplicated-sequence views; the missing sequence is not assigned an
invented rank, and repeated entries are not counted as independent proteins.

We distinguish three evidence tiers. The conservative primary-paper set
contains 12 unique sequences corresponding to 15 entries, with source-level
activity evidence rechecked in the existing audit. Including patent evidence
expands this set to 24 sequences and 36 entries. The broader workbook-active
set contains 81 sequences and 102 entries; the remaining 42 unique sequences
have only conditional non-detection evidence. Not every broader workbook
assay has been independently verified, and non-detection under one assay
condition is not a claim of universal inactivity. The workbook's curated
``Ranked Candidates'' list is not used as binary experimental ground truth.

\subsubsection{Deployment protocol and evaluation}

The primary deployment uses the Reaction-Sim-trained checkpoint, with its
recorded recipe frozen before that version's Case1 scoring. No Case1 labels
are used to fit the retrieval parameters. All other target-trained checkpoints
are reported as transfer diagnostics without replacing the primary checkpoint
on the basis of Case1 performance. Because the reaction and earlier outcomes
were examined during development, this is retrospective evidence rather than
an untouched external validation cohort.

For each candidate view, we report primary-paper recovery at ranks 5, 10,
and 25 and normalize recall by that view's own positive count. The unique
view asks how many distinct documented sequences would be selected; the
entry view preserves the requested workbook panel, including duplicate
sequences and its deterministic tie order. As a descriptive reference, a
uniform random ranking has expected primary-paper recovery at 25 of
$25(12/123)=2.439$ unique sequences or $25(15/144)=2.604$ entries. These are
expectations, not uncertainty estimates for the model.

We additionally report conditional AUROC for the 81 workbook-active versus
42 non-detect-only unique sequences. The same enzyme can appear in several
constructs or studies, and assay conditions differ across sources. We
therefore do not attach candidate-IID confidence intervals or interpret this
AUROC as prospective activity-prediction accuracy.

\subsubsection{Results and transfer limits}

Table~\ref{tab:exp-case1} and Figure~\ref{fig:exp-case1} show the primary
Reaction-Sim results. Original \methodname\ recovers 10/12 unique
paper-supported catalysts at rank 25, compared with 9/12 for its native F3
encoder before phase 2. Both recover 10/15 paper-supported entries in the
144-entry ranking. With phase-2 biological supervision, unique recovery
increases to 11/12, whereas entry recovery remains 10/15. Conditional AUROC
changes from 0.612875 to 0.613757. The unique-sequence gain is one catalyst,
TsT4Ease WT (H008), moving from rank 26 to rank 25. Cofactor removal retains
this gain; shuffled labels and mechanism removal do not.

\begin{table}[htbp]
\centering\small
\caption{Case1 primary deployment and matched phase-2 controls. Unique
recovery is out of 12 positives among 123 sequences; entry recovery is out
of 15 positives among 144 resolved entries from the original 145-entry
catalogue. AUROC is conditional on the broader 81/42 workbook labels.
Controls use the same Reaction-Sim F3, phase-2 budget, and inference settings.
``Native F3'' is the saved pre-phase-2 diagnostic.}
\label{tab:exp-case1}
\resizebox{\linewidth}{!}{% Generated from evidence_snapshot.json by prepare_tables.py.
\begin{tabular}{lrrrrr}
\toprule
Model / labels & Unique @5 & Unique @10 & Unique @25 & Entries @25 & AUROC \\
\midrule
Native F3 & 1 & 3 & 9 & 10 & 0.600823 \\
\methodname & 0 & 3 & 10 & 10 & 0.612875 \\
\vfourBio & 0 & 3 & 11 & 10 & 0.613757 \\
Shuffled labels & 0 & 3 & 10 & 10 & 0.611111 \\
Without EC & 0 & 3 & 10 & 10 & 0.613757 \\
Without cofactor & 0 & 3 & 11 & 10 & 0.614345 \\
Without mechanism & 0 & 3 & 10 & 10 & 0.613169 \\
\bottomrule
\end{tabular}
}
\end{table}

\begin{figure}[htbp]
\centering
\includegraphics[width=\linewidth]{\expFigPath 08_case1_controls.pdf}
\caption{Case1 biological-supervision controls. Left: recovery at rank 25
among 123 unique sequences with 12 paper-supported catalysts. Middle:
recovery among 144 resolved entries with 15 paper-supported entries. Right:
conditional AUROC differences from original \methodname\ (0.612875).
All variants use the Reaction-Sim lineage and biological weight 1. The
one-sequence gain in the deduplicated view does not improve the entry-level
count. Colors, hatching, and printed counts distinguish the two evaluation
units. These are retrospective literature results, not new wet-lab assays.}
\label{fig:exp-case1}
\end{figure}

Early cutoffs remain challenging: the primary final model recovers no
primary-paper catalyst in its first five unique sequences and three in its
first ten. Transfer also depends on the source training dataset: the
EnzymeMap-trained \methodname\ checkpoint recovers 6/12 at rank 25 with
conditional AUROC 0.470606 (Appendix~\ref{app:exp-case-targets}). In the
existing homology audit, four of the ten primary catalysts recovered by the
Reaction-Sim model have retrieved training homologs with at least 90\%
identity; both query molecules also co-occur within larger training
participant sets. Missing qualifying homology hits do not prove sequence
novelty, and this audit does not cover frozen-backbone pretraining exposure.
\FloatBarrier

\subsection{Ablation and sensitivity analyses}
\label{sec:exp-ablation}

We retain the recorded controls that isolate a specific training or inference
choice. Inference comparisons use identical trained weights and candidate
pools. Phase-2 label controls hold F3 fixed. The replicated temperature and
fresh-F3 biological studies use their own matched recipes and seeds; they
are not presented as independent replications of the final single-seed V4
result. Appendix~\ref{app:exp-coverage} lists the remaining component controls
needed for a complete architectural ablation.

\subsubsection{Dictionary weight and residual strength}

At residual multiplier $c=1$, increasing dictionary weight from 0.25 to 0.40
raises Reaction-Sim \expEtoR\ MRR from 0.523804 to 0.537749. A further increase
to 0.50 reduces \expRtoE\ MRR on all three ReactZyme splits, while increasing
EnzymeMap EF5 and EF10 in both pools. At $\alpha=0.40$, reducing $c$ from 1
to 0.5 lowers BEDROC$_{85}$ in both screening settings. These effects show
a trade-off between retrieval directions and screening budgets, rather than
a uniformly beneficial inference adjustment
(Figures~\ref{fig:exp-inference-rz} and~\ref{fig:exp-inference-em}).

\begin{figure}[htbp]
\centering\includegraphics[width=\linewidth]{\expFigPath 01_reactzyme_inference.pdf}
\caption{ReactZyme inference sensitivity with fixed checkpoints. Columns
change one coefficient at a time; rows retain both directions and all three
splits. Values are signed differences in MRR multiplied by 100; for example,
$+1.3945$ represents $+0.013945$ raw MRR. The color scale is linear, centered
at zero, with blue for positive and orange for negative changes. Larger
Enzyme-Sim/Time \expEtoR\ effects require the near-tie check in
Figure~\ref{fig:exp-ties}.}
\label{fig:exp-inference-rz}
\end{figure}

\begin{figure}[htbp]
\centering\includegraphics[width=\linewidth]{\expFigPath 02_enzymemap_inference.pdf}
\caption{EnzymeMap inference sensitivity. Blue circles denote the full pool
(Table 1 protocol); open orange squares denote training-ID exclusion (Table
2 protocol). Each point is the changed setting minus its stated reference;
positive values favor the change. BEDROC differences are in raw units,
not percentages. The four metric-specific axes prevent conflating enrichment
with very-early recognition. Single-seed checkpoints and candidate pools are
held fixed.}
\label{fig:exp-inference-em}
\end{figure}
\FloatBarrier

\subsubsection{Anchor-balanced training and inverse temperature}

The available replicated comparison isolates inverse temperature within an
earlier 18-epoch EnzymeMap recipe: $\beta=5$ versus $\beta=10$, with paired
seeds 17, 42, and 73 and the remaining training and phase-2 choices fixed.
It does not isolate anchor balancing against the original MLNCE loss.
Lower inverse temperature improves all BEDROC$_{85}$, BEDROC$_{20}$, and EF5
seed pairs. Overall, 22 of the 24 seed-by-metric comparisons improve; seed
73 produces the two negative EF10 differences. Full-pool mean BEDROC$_{85}$
increases from $0.494040\pm0.035007$ to $0.521656\pm0.037549$, where the
dispersion is sample standard deviation over three seeds. Figure~\ref{fig:exp-temp}
retains every seed rather than displaying only the mean. The metric cells
are correlated, so their count is not a count of independent experiments.

\begin{landscape}
\begin{figure}[p]
\centering\includegraphics[width=\linewidth]{\expFigPath 03_temperature_paired_seeds.pdf}
\caption{Paired inverse-temperature effects on EnzymeMap at fixed 18 epochs.
Values are $\beta=5$ minus $\beta=10$; circles are individual seeds and dark
diamonds their arithmetic mean. Upper row: full pool (Table 1). Lower row:
training-ID-excluded pool (Table 2). Scales match across rows for each metric;
BEDROC is unscaled. These observations describe a matched component study,
not three-seed replication of the final 10-epoch V4 recipe or confidence
intervals around that recipe.}
\label{fig:exp-temp}
\end{figure}
\end{landscape}

\subsubsection{Biological supervision in phase 2}

We compare the unannotated V4 control with correctly assigned labels,
shuffled labels, and removal of EC, cofactor, or transformation supervision.
All active family coefficients are 1, and the averaging divisor remains
three when a family is removed. Thus removal tests the family contribution
at fixed remaining coefficients and also reduces total regularization;
it is not an equal-total-weight reallocation.

The controls do not establish an overall benchmark advantage from adding
biology. Full-pool BEDROC$_{85}$ is 0.572337 without added supervision,
0.572096 with correct labels, and 0.572093 with shuffled labels. Every
displayed phase-2 variant slightly lowers both BEDROC metrics in both pools,
while EF effects are mixed. ReactZyme point differences are likewise mixed
(Figures~\ref{fig:exp-bio-rz} and~\ref{fig:exp-bio-em}). The Case1 controls
in Figure~\ref{fig:exp-case1} show a conditional top-25 crossing, not a
consistent improvement across evaluation units and benchmarks.

\begin{figure}[htbp]
\centering\includegraphics[width=\linewidth]{\expFigPath 04_reactzyme_phase2_biology.pdf}
\caption{ReactZyme phase-2 biological-supervision controls, each minus
original \methodname. Cell values are $100\Delta\mathrm{MRR}$. Colors use
one shared symmetric logarithmic scale, linear within $\pm0.001$ displayed
units ($\pm10^{-5}$ raw MRR), so larger Time \expEtoR\ differences do not
wash out small effects elsewhere. Printed values are not log transformed;
color intensity is nonlinear and is not directly comparable with the linear
scale in Figure~\ref{fig:exp-inference-rz}. Correct/shuffled labels use
weight 1; removal controls retain other coefficients and the divisor of three.
Visible color does not imply a practically large effect.}
\label{fig:exp-bio-rz}
\end{figure}
\FloatBarrier

\begin{landscape}
\begin{figure}[p]
\centering\includegraphics[width=\linewidth]{\expFigPath 05_enzymemap_phase2_biology.pdf}
\caption{EnzymeMap phase-2 biological-supervision controls. Each variant is
compared with original V4 without added supervision; F3, architecture,
inference, pools, and training budget are fixed. Upper/lower rows are the
full/training-ID-excluded pools. Filled blue circles denote correct labels,
open squares shuffled labels, and gray circles family-removal controls.
Metric scales match across rows. Axis multipliers matter: the BEDROC changes
are on the order of $10^{-4}$ in raw units. Correct labels and shuffled labels
nearly coincide on BEDROC$_{85}$.}
\label{fig:exp-bio-em}
\end{figure}
\end{landscape}

\subsubsection{Near-tie sensitivity of reaction retrieval}
\label{sec:exp-ties}

Some ReactZyme \expEtoR\ metrics are sensitive to the ordering of nearly
equal scores. We therefore retain the archived optimistic and pessimistic
rank bounds obtained by permitting alternative orderings within a
$10^{-6}$ score band. These are numerical sensitivity bounds, not confidence
intervals. The small Reaction-Sim difference remains distinct in this
diagnostic, whereas Enzyme-Sim and Time bounds overlap substantially
(Figure~\ref{fig:exp-ties}). For example, the Time point estimate increases
from 0.780669 to 0.800614 under phase-2 biology, but that increase alone
does not establish robust improvement in reaction discrimination.

\begin{figure}[htbp]
\centering\includegraphics[width=\linewidth]{\expFigPath 06_reactzyme_near_tie_sensitivity.pdf}
\caption{Recorded \expEtoR\ MRR and numerical near-tie bounds for original
V4 and phase-2 weight-1 biology. Points are recorded scores; line segments
span optimistic and pessimistic ranks under a $10^{-6}$ score band. Each
split has its own zoomed horizontal scale. Segment lengths should be compared
within a panel, not across panels. The intervals quantify ranking sensitivity
and have no sampling-uncertainty interpretation.}
\label{fig:exp-ties}
\end{figure}
\FloatBarrier

\subsubsection{Biological supervision during F3 training}

A separate EnzymeMap study retrains F3 with relative biological-loss weights
0, 0.1, or 1 across paired seeds 42, 43, and 44, keeping phase 2 unannotated.
This directly changes F3 training, unlike \vfourBio. Weight 1 increases
mean full-pool BEDROC$_{85}$ by 0.010256 but decreases mean EF5 by 0.151444.
Only two of three seeds improve BEDROC$_{85}$; Table-2 EF5 decreases in all
three seeds. Weight 0.1 improves Table-2 EF10 in all three seeds, but not all
other metrics. Figure~\ref{fig:exp-f3-bio} exposes this seed and objective
dependence. These results support a mixed biological-loss trade-off rather
than a reproducible improvement on every early-screening endpoint.

\subsection{Interpretation and remaining validation}

The recorded experiments support improved EnzymeMap screening over the
released CLIPZyme checkpoint and the displayed FGW-CLIP point estimates,
and competitive ReactZyme scores relative to published main-method results.
They also expose three boundaries: a newly retrained parent-model comparator
is stronger on the six ReactZyme MRR cells; biological supervision has mixed
or negligible controlled effects; and Case1 remains one dependent,
retrospective literature panel. A claim of general superiority or broad
real-world generalization would require the remaining matched controls,
replication of the final recipe across seeds, and a prospectively specified
external reaction panel. These unresolved questions are retained alongside
the favorable benchmark results.

\begin{landscape}
\begin{figure}[p]
\centering\includegraphics[width=\linewidth]{\expFigPath 07_f3_biology_paired_seeds.pdf}
\caption{Fresh-F3 biological-loss study on EnzymeMap. Points are each
biological variant minus its same-seed unannotated control; the last row in
each panel is the arithmetic mean. Blue circles and open orange squares
indicate weights 0.1 and 1. Upper/lower panels denote the full/training-ID-excluded
pools, with matching metric scales. Phase 2 stays unannotated, and seed 42
uses the archived control. All seeds are displayed; these are not
confidence intervals or estimates from the phase-2-only \vfourBio\ model.}
\label{fig:exp-f3-bio}
\end{figure}
\end{landscape}

